\capituloPregunta{¿Cómo modelar una respuesta continua?}

\section{Motivación}

Comparar tratamientos permite decidir si existen diferencias. Modelar permite preguntar cuánto cambia una respuesta cuando se modifica una condición.

\section{Intuición}

Un modelo no es la realidad. Es una representación simplificada que debe revisarse. Una recta puede describir una tendencia, pero también puede ocultar curvatura, valores atípicos o errores de medición.

\section{Elementos mínimos}

\begin{itemize}
  \item variable respuesta;
  \item variable explicativa;
  \item rango de observación;
  \item residuos;
  \item interpretación de pendiente;
  \item límites de extrapolación.
\end{itemize}

\section{Actividad}

Seleccione una variable experimental continua, por ejemplo intensidad de luz, altura, tiempo o temperatura. Proponga un modelo lineal e indique qué gráfica de residuos revisaría antes de confiar en la predicción.

\begin{cuidado}
No extrapole fuera del rango observado sin justificación. Un modelo puede ajustar bien dentro del experimento y fallar fuera de él.
\end{cuidado}

\section{Conexión}

A veces la comparación de tratamientos debe ajustarse por variables adicionales. Ese es el punto de entrada a ANCOVA.
